% !TeX spellcheck = it_IT

\documentclass[format=169, handout]{beamer}

\usepackage[utf8]{inputenc}
\usepackage[italian]{babel}
\usepackage{hyperref}
\usepackage{xcolor}
\usepackage{listings}
\usepackage{forest}


% for matlab code
\definecolor{mygreen}{RGB}{28,172,0} % color values Red, Green, Blue
\definecolor{mylilas}{RGB}{170,55,241}
\definecolor{mGray}{rgb}{0.5,0.5,0.5}
\definecolor{mPurple}{rgb}{0.58,0,0.82}
\definecolor{backgroundColour}{rgb}{0.95,0.95,0.92}

\lstdefinestyle{matlab}{language=Matlab,%
	%basicstyle=\color{red},
	breaklines=true,%
	morekeywords={matlab2tikz},
	keywordstyle=\color{blue},%
	morekeywords=[2]{1}, keywordstyle=[2]{\color{black}},
	identifierstyle=\color{black},%
	stringstyle=\color{mylilas},
	commentstyle=\color{mygreen},%
	showstringspaces=false,%without this there will be a symbol in the places where there is a space
	numbers=left,%
	numberstyle={\tiny \color{black}},% size of the numbers
	numbersep=9pt, % this defines how far the numbers are from the text
	emph=[1]{for,end,break},emphstyle=[1]\color{red}, %some words to emphasise
	%emph=[2]{word1,word2}, emphstyle=[2]{style},    
}

\usetheme{metropolis}

\usepackage{mathtools}
\DeclarePairedDelimiter{\ceil}{\lceil}{\rceil}

\title{Esercitazioni di Informatica B}
\subtitle{Diagrammi}

\author{Stefano Cereda\\
	stefano.cereda@polimi.it
}
\date{11/12/2018}
\institute[PoliMi]{\vspace{0.5cm}\centering Politecnico di Milano \\ \vspace{0.2cm}
	\includegraphics[width=0.2\textwidth]{../logopolimi}}

\setbeamercovered{invisible}

\makeindex

\begin{document}
	\begin{frame}
	\maketitle
\end{frame}

\begin{frame}[allowframebreaks]{Caldaia}
	In un cinema di Milano la temperatura nella sala è regolata in modo
	automatico; un sensore monitora la temperatura rilevando un valore ogni
	minuto. I valori letti vengono inseriti in un array temp.
	
	La caldaia si accende solo quando la temperatura rilevata é inferiore ad una
	certa soglia e si spegne non appena viene raggiunta tale soglia.
	\begin{enumerate}
		\item  Scrivere in Matlab uno script che:
		\begin{enumerate}
			\item Generi in maniera random un vettore di temperature e che chieda
			in ingresso un valore soglia ed un parametro ``costoAlMinuto''
			indicante il costo al minuto del gas consumato dalla caldaia;
			\item Calcoli e stampi il costo totale del gas consumato ed un vettore
			contenenti gli indici dell’array temp corrispondenti ai minuti nei
			quali la caldaia é rimasta accesa.
			
			Per il calcolo del costo totale del gas consumato si consideri che se il
			totale dei minuti nei quali la caldaia ha funzionato supera i 30 minuti,
			allora il costo al minuto va diminuito del 20\%.
		\end{enumerate}
		\item Disegnare il grafico della temperatura al variare del tempo, evidenziando
		sul grafico stesso:
		\begin{enumerate}
			\item Il valore soglia con una retta orizzontale nera a doppio spessore
			\item I minuti nei quali la caldaia ha funzionato, marcandoli con il simbolo
			‘X’ rosso (senza congiungerli).
		\end{enumerate}
	\end{enumerate}
\end{frame} 

\begin{frame}[allowframebreaks, fragile]{Roulette}
	Nel gioco della roulette si ha una ruota raffigurante 37
	numeri da 0 a 36. Ad ogni turno un giocatore può
	scommettere su un numero o su particolari combinazioni
	degli stessi. Alla fine delle puntate la ruota selezionerà il
	numero vincente.
	
	\begin{enumerate}
		\item Si scriva una funzione che simuli la ruota e generi il
		numero vincente utilizzando la funzione \emph{rand}.
		\item Si scriva una funzione che provi ad eseguire N giri e salvi
		i risultati delle frequenze di ogni numero in un array di
		37 elementi (il primo elemento per lo zero, il secondo per
		l'uno e così via).
		\item Si tracci il grafico delle frequenze prima utilizzando
		\emph{rand} come generatore e poi utilizzando la funzione
		\emph{rand\_truccata} mostrata di seguito e commentare il
		risultato.
		\item Qual è la probabilità di vincere se si scommette
		sull'uscita di un numero dispari? Scrivere una funzione
		che calcoli questa probabilità effettuando 100000
		simulazioni di giocate.
	\end{enumerate}
	
	\framebreak
\begin{lstlisting}[style=matlab]
function [n]=rand_truccata()
r=rand();
if r<0.02	
    n=0;
else
    n=rand();
end
\end{lstlisting}
\end{frame}

\begin{frame}{Dado}
	Nel gioco dei dadi ad ogni faccia del dado è associato un
	numero da 1 a 6.
	
	\begin{enumerate}
		\item Si scriva una funzione che simuli il dado e generi il
		numero vincente utilizzando la funzione rand.
		\item Si scriva una funzione che provi ad eseguire N lanci e
		salvi i risultati delle frequenze di ogni numero in un array
		di 6 elementi (il primo elemento per l’uno, il secondo per
		il due e così via).
		\item Qual è la probabilità di vincere se si scommette
		sull'uscita di un numero? Scrivere una funzione che
		calcoli questa probabilità effettuando 100000 simulazioni
		di lanci.
	\end{enumerate}
\end{frame}

\iffalse
\begin{frame}[fragile]{Paraboloide}
	Si disegni la superficie del paraboloide iperbolico dato dalla seguente
	equazione:
\begin{lstlisting}[style=matlab]
z=x.^2-y.^2
\end{lstlisting}

	per valori di x e y nell'intervallo [-5,+5].\\
	Impostare titoli ed etichette.
\end{frame}
\fi


\end{document}